\documentclass[10pt,a4paper]{amsart}
\usepackage[margin=19mm]{geometry}
\usepackage{amsmath,amssymb,mathtools}
\usepackage[scr=rsfs,cal=euler]{mathalfa}
\usepackage{xfrac}
\usepackage[unicode,hidelinks,bookmarks=false]{hyperref}
\newcommand{\ie}{i.e.\ }
\newcommand{\cf}{cf.\ }
\newcommand{\resp}{respectively\ }
\newcommand{\Subsection}{\subsection}
\providecommand{\nobreakdash}{\nobreak\hbox{-}}
\setlength{\emergencystretch}{2em}
\multlinegap0pt
\usepackage{bm}
\usepackage{bbm}
\usepackage{dsfont}
\usepackage{xspace}
\usepackage{cancel}
\usepackage{mathtools}
\usepackage{amsthm}
\makeatletter
\@ifundefined{theorem}{\newtheorem{theorem}{Theorem}}{}
\@ifundefined{fact}{\newtheorem{fact}[theorem]{Fact}}{}
\@ifundefined{proposition}{\newtheorem{proposition}[theorem]{Proposition}}{}
\makeatother
\title[A step toward Chen-Lih-Wu conjecture]{A step toward Chen-Lih-Wu conjecture}
\author{Yangyang Cheng, Zhenyu Li, Wanting Sun, Guanghui Wang}
\date{Source: arXiv:2511.03957v1 (CC BY 4.0)}
\begin{document}
\maketitle

\noindent\emph{Source and licence.} A step toward Chen-Lih-Wu conjecture. Version arXiv:2511.03957v1 (CC BY 4.0), distributed under the Creative Commons Attribution 4.0 International licence, as recorded in the arXiv licence metadata for this exact version and shown on the abstract page https://arxiv.org/abs/2511.03957v1. Full rights notice: licenses/arxiv-2511.03957-CC-BY-4.0.txt.\par\smallskip

\noindent\textit{Editorial scope.} Counted unit: the Proposition whose source label is \texttt{PMatching} in the article (source0.tex lines 641--654, source label \texttt{PMatching}), delivered together with the article's own complete original proof of it (source0.tex lines 656--838, one \texttt{proof} environment of 183 lines). No mathematics is written, repaired or re-derived: statement and proof are the article's own text line for line. Required context retained uncounted: MMProperty (lines 606--609). Every definition, display and label of those lines is retained unchanged. Omitted from this package and not redistributed: 1--605, 610--640, 655--655, 839--4122. Nothing inside a retained excerpt is omitted or altered: every retained body is delivered exactly as the article's lines are written, so no omission receipt of any kind accompanies this package. Citations: the retained excerpts carry no citation command at all, so no bibliography entry of the article is transcribed and no reference is redistributed. Reference adaptation: none was needed and none was made; every reference inside the delivered unit resolves inside it, so no unresolved reference or citation is produced. Printed slips kept literal: no printed word, symbol, hypothesis or number of the counted statement or of its proof was altered. Layout and class: the article's own class is \texttt{article} with packages including t1enc, inputenc, babel, amsmath, amsthm, mathrsfs, amsfonts, latexsym, graphicx, tikz, float, xcolor, algorithm, algorithmic; the wrapper is the lane's pinned \texttt{amsart} editorial wrapper at 10pt on A4, which restates the theorem-like environments the retained text needs and adds the article's own macro definitions that the retained text needs (none), with extra wrapper packages (bm, bbm, dsfont, xspace, cancel, mathtools). The delivered numbering is the unit's own. Assets: no retained span contains a figure environment, a graphics-inclusion command, a PDF-page inclusion, a table or a TikZ picture; no image file is copied into this package, the package declares no asset dependency and no image file is redistributed with it. Licence: version arXiv:2511.03957v1 of this article is distributed under the Creative Commons Attribution 4.0 International licence, as recorded in the arXiv licence metadata for this exact version and shown on the abstract page https://arxiv.org/abs/2511.03957v1. A full rights notice is delivered at licenses/arxiv-2511.03957-CC-BY-4.0.txt. Characters: every retained excerpt is pure ASCII, so the delivered engine, XeLaTeX with its default Latin Modern text font, reports no missing character.
\begin{fact}\label{MMProperty}
   Let $M$ be a maximum matching of a graph $G$ and $x,y$ be two distinct vertices in $V(G)\setminus V(M)$. %For any vertex $v\in V(G)\setminus V(M)$, define $SN(v):=\{z\in V(M):wz\in M \text{ for some } w\in N(v)\}$. 
   Then $E(G[SN(x),SN(y)])=\emptyset$. %({\color{red}define SN(v) in Notation?}) %  If $M$ is the maximum matching of $G$, then, for any two different vertices $x,y\in V(G)\setminus V(M)$, we have that $E(SN(x),SN(y))=\emptyset$, where $SN(v):=\{z\in V(M):wz\in M \text{ for some } w\in N(v)\}$ for any $v\in V(G)\setminus V(M)$.
\end{fact}

\begin{proposition}\label{PMatching}
Let $n$ be a positive even integer and $\gamma$ be a constant such that $0<{1}/{n}\ll\gamma\ll 1$. Suppose that $G$ is a graph on $n$ vertices so that 
\begin{align}\label{B2Ore}
\sigma(G)\ge n-\gamma n.
\end{align}
Then one of the following holds:
\begin{itemize}
\setlength{\parskip}{0pt}
    \item[{\rm (i)}] $G$ admits a perfect matching;
    \item[{\rm (ii)}] $G$ contains a $2\gamma$-independent set of size at least $\frac{n}{2}$;
    \item[{\rm (iii)}] $G$ consists of two odd components $C_1$ and $C_2$. Moreover, if $|C_i|\leq \frac{1-\gamma}{2}n$ for some $i\in [2]$, then % the size of a component $C$ is at most $\frac{1-\gamma}{2}n$ then 
    $C_i$ is a clique.
\end{itemize}
\end{proposition}

\begin{proof}
Suppose that $G$ has no perfect matching. Let $M$ be a maximum matching in $G$. Then, $|V(M)|\leq n-2$ and there are two  distinct vertices $x,y\in V(G)\setminus V(M)$. %Define $SN(v):=\{z'\in V(M):zz'\in E(M) \text{ for some } z\in N(v)\}$ for each $v\in V(G)\setminus V(M)$.  
The maximality of $M$ implies that $N(x)\cup  N(y)\subseteq  V(M)$, and so $SN(x)\cup SN(y)\subseteq V(M)$. Since $xy\notin E(G)$,  % For each vertex $x\notin V(M)$, we have that $d(x)=|N(x)|=|SN(x)|$. 
by \eqref{B2Ore} we have 
\begin{align}\label{SNOre}
   d(x)+d(y)= |SN(x)|+|SN(y)|\ge n-\gamma n.
\end{align}
% {\color{red}(delete??)Fact \ref{MMProperty} implies that $E(SN(x),SN(y))=\emptyset$, which also implies $N(x)\cap SN(y)=\emptyset$.  %Since $SN(x)\cup SN(y)\subseteq V(M)$, we have
% Therefore, 
% $$|E(M)|\ge |SN(x)\cap SN(y)|+\frac{|SN(x)\Delta SN(y)|}{2}\ge \frac{n-\gamma n}{2}.$$}
Define $SN(x,y):=SN(x)\cap SN(y)$. Fact \ref{MMProperty} implies that $SN(x,y)$ is an independent set of $G$ and  
\begin{align}\label{degreesum}
    d(x)+d(y)\leq 2e(M)\leq n-2.
\end{align}
We divide our proof into two cases.

\medskip
{\bf Case 1. } 
$SN(x,y)\ne \emptyset$. 
\medskip

In this case, we will show that (ii) holds. 
Without loss of generality, assume that $d(x)\le d(y)$. Then \eqref{degreesum} implies $d(x)<\frac{n}{2}$. Choose $z\in SN(x,y)$ and assume $zz'\in E(M)$. Hence  $z'\in N(x)\cap N(y)$ and  $N(z)\subseteq V(M)$.  Thus, \eqref{B2Ore} implies that 
$$
d(z)\ge n-\gamma n-d(x)>\frac{n}{2}-\gamma n.
$$
Notice that $N(z)\cap (SN(x)\cup SN(y))=\emptyset$. Otherwise, there exists a matching of size $|E(M)|+1$, a contradiction. Thus, 
$|SN(x)\cup SN(y)|\le n-d(z)<\frac{n}{2}+\gamma n$. Together  with \eqref{SNOre}, one has  
$$|SN(x)\cap SN(y)|= |SN(x)|+|SN(y)|-|SN(x)\cup SN(y)|\ge \frac{n}{2}-2\gamma n.$$
Recall that $SN(x,y)$ is an independent set in $G$. By adding at most $2\gamma n$ arbitrary vertices to $SN(x,y)$, we obtain a $2\gamma$-independent set of size at least $\frac{n}{2}$ in $G$, as desired. 



%In this case, we will show that (ii) holds. % $G$ contains a $3\gamma$-independent set of size at least $n/2$. 
%Since $M$ is not a perfect matching, one has $|E(M)|\le \frac{n}{2}-1$. Fact \ref{MMProperty} implies that $N(x)\cap SN(y)=\emptyset$. Moreover, $N(x),SN(y)\subseteq V(M)$. Thus, $|N(x)|+ |SN(y)|\le n-2$. 
%Without loss of generality, assume that $d(x)\le d(y)$. Choose $z\in SN(x,y)$ and assume $zz'\in E(M)$. Hence  $z'\in N(x)\cap N(y)$ and  $N(z)\subseteq V(M)$.  Thus, \eqref{B2Ore} implies that 
%$$
%d(z)\ge n-\gamma n-d(x).
%$$

%If $d(x)\le \frac{n}{2}-2\gamma n$, then $d(z)\geq \frac{n}{2}+\gamma n$. Similarly, \eqref{SNOre} implies $d(y)\ge \frac{n}{2}+\gamma n$. 
%Thus, $d(y)+d(z)\geq n+2\gamma n$, that is, $(N(z)\cap SN(y))\setminus \{z'\}\ne \emptyset$. Assume $w\in (N(z)\cap SN(y))\setminus \{z'\}$ and $ww'\in E(M)$. Then $M':=(M \setminus\{zz',ww'\})\cup \{xz', yw',zw\}$ is a matching of size $|E(M)|+1$, a contradiction. 

%If $d(x)> \frac{n}{2}-2\gamma n$, then by \eqref{degreesum} one has  $d(x)\le \frac{n}{2}\leq d(y)\le \frac{n}{2}+2\gamma n$. Therefore, $d(z)\geq \frac{n}{2}-3\gamma n$. % Consider vertices $x,z$, \eqref{B2Ore} implies that $$d(z)\ge n-\gamma n -d(x) > \frac{n}{2}-\gamma n.$$ 
%Notice that $N(z)\cap (SN(x)\cup SN(y))=\emptyset$. Thus, 
%$|SN(x)\cup SN(y)|\le |V(M)|-d(z)\le \frac{n}{2}+3\gamma n-2$. Together  with \eqref{SNOre}, one has  
%$$|SN(x)\cap SN(y)|= |SN(x)|+|SN(y)|-|SN(x)\cup SN(y)|\ge \frac{n}{2}-4\gamma n.$$
%Recall that $SN(x,y)$ is an independent set in $G$. By adding at most $4\gamma n$ arbitrary vertices to $SN(x,y)$, we obtain a $5\gamma$-independent set of size at least $\frac{n}{2}$ in $G$, as desired. 

\medskip
{\bf Case 2.} %For each maximum matching $M$ and any two vertices $x,y\in V(G)\setminus V(M)$, we have 
$SN(x,y)=\emptyset.$
\medskip

In this case, one may assume that 
$SN(x,y)=\emptyset$ for every maximum matching $M$ of $G$ and any two distinct  vertices $x,y\in V(G)\setminus V(M)$. Suppose that  $|V(G)\setminus V(M)|>2$. Then there is a vertex $z\in V(G)\setminus (V(M)\cup \{x,y\})$. It follows from  \eqref{B2Ore} that 
$$d(x)+d(y)+d(z)\ge \frac{3}{2}(n-\gamma n)>|V(M)|.$$ 
Recall that $N(x)\cup N(y)\cup N(z)\subseteq V(M)$. Therefore, there are two vertices $u,v\in \{x,y,z\}$ such that $N(u)\cap N(v)\neq \emptyset$, thus $SN(u,v)\neq \emptyset$, a contradiction. Hence, $|V(G)\setminus V(M)|=2$. We will show that (iii) holds. 

%, Fact \ref{MMProperty} implies that, for any three different vertices $x_1,x_2,x_3\in  V(G)\setminus V(M)$, we have $$e(G[SN(x_i),SN(x_j)])=0,$$
%where $i,j\in \{1,2,3\}$ and $i\ne j$. \eqref{B2Ore} implies that $$d(x_1)+d(x_2)+d(x_3)\ge \frac{3}{2}(n-\gamma n)>|V(M)|$$, that is, there exists $x_i\ne x_j$ such that $SN(x_i)\cap SN(x_j)\ne \emptyset$, contradicting the assumption.

%If $|V(G)\setminus V(M)|=2$, then d
%Define $R:=V(G)\setminus (SN(x)\cup SN(y)\cup\{x,y\})$. %In fact, one may assume that the maximum matching $M$and vertices $\{x,y\}$ are chosen such that $|R|$ is {\color{red} minimum}. 
%We first consider $R=\emptyset$. Recall that Fact~\ref{MMProperty} implies  $E(G[SN(x),SN(y)])=\emptyset.$ Together with $SN(x,y)=\emptyset$, one may partition the edges in $M$ into two parts, say $M_x\cup M_y$, where $M_x$ denotes the set of edges in $M$ whose both endpoints are adjacent to $x$ and $M_y$ is defined similarly. %It suffices to show that $N(x)\cap SN(y)=\emptyset$ and $N(y)\cap SN(x)=\emptyset$. Indeed, we only need to show the former one. Assume $v\in N(x)\cap SN(y)$, and $vv'\in E(M)$. Then $v'\in SN(x)$, contradicting Fact \ref{MMProperty}.
% If $R=\emptyset$, then by the fact that $SN(x,y)=\emptyset$ 
% and Fact \ref{MMProperty},  %Notice that by replacing  $x$ with any vertex $z\in V(M_x)$, 
% Notice that $G[V(M_x),V(M_y)]$ is an empty graph. Otherwise, suppose that there exists an edge $uv$ such that $u\in V(M_x)$ and $v\in V(M_y)$. Let $uu',vv'\in E(M)$. Hence,  replacing $\{uu',vv'\}$ to $\{xu',uv,v'y\}$ in $M$ obtains a larger matching $M'$, contradicting that matching $M$ is maximum. 
%Therefore, $G[V(M_x)\cup \{x\}]=:C_1$ and $G[V(M_y)\cup \{y\}]=:C_2$ are two connected odd components of $G$, as desired. % forms a clique with odd order. Similarly, $G[V(M_y)\cup \{y\}]$ also forms a clique with odd order, as desired.

%If there exists an edge $uv$ such that $u\in V(M_x)\cup \{x\}$ and $v\in V(M_y)\cup\{y\}$, then replacing $\{uu',vv'\}$ to $\{xu',uv,v'y\}$ in $M$ obtains a larger matching $M'$, contradicting that matching $M$ is maximum.

%{\color{red}By symmetry and the maximality of $|R|$, we obtain that for any vertex $z\in V(M_x)$, it is adjacent to all vertices in $(V(M_x)\cup \{x\})\setminus \{z\}$. Therefore, $G[V(M_x)\cup \{x\}]$ forms a clique with odd order. Similarly, $G[V(M_y)\cup \{y\}]$ also forms a clique with odd order, as desired.}

%Now, we consider $R\neq \emptyset$. 
%Define $R:=V(G)\setminus (SN(x)\cup SN(y)\cup\{x,y\})$. 
Since $SN(x,y)=\emptyset$ and $E(G[SN(x),SN(y)])=\emptyset$, there are two vertex-disjoint subgraphs $G_1:=G[\{x\}\cup SN(x)]$ and $G_2:=G[\{y\}\cup SN(y)]$ in $G$ such that $E(G[V(G_1),V(G_2)])=\emptyset$. Define $R:=V(G)\setminus (V(G_1)\cup V(G_2))$. If $R=\emptyset$, then the two endpoints of each edge in $M$ are adjacent to the same vertex in $\{x, y\}$. It follows that $G$ consists of two connected odd components $G_1$ and $G_2$, as desired. Suppose that $R\neq \emptyset$. In order to prove (iii) holds, it is necessary to assign vertices in $R$ to $G_1$ or $G_2$. Choose a vertex $v \in R$ with $vv'\in E(M)$. Hence $xv',yv'\notin E(G)$. Together with $xy\notin E(G)$ and \eqref{B2Ore}, we have 
%\begin{align*}
%    &d(x)+d(y)=|SN(x)|+|SN(y)|\geq n-\gamma n;\\
%    &d(x)+d(v')=|SN(x)|+|N(v')|\geq n-\gamma n;\\
%    &d(y)+d(v')=|SN(y)|+|N(v')|\geq n-\gamma n.
%\end{align*}
%Therefore,
$$
d(x)+d(y)+d(v')=|SN(x)|+|SN(y)|+|N(v')|\geq \frac{3}{2}(n-\gamma n).
$$
Recall that $SN(x,y)=\emptyset$. Hence, 
\begin{align}\label{eitheror}
    \text{either}\ SN(x)\cap N(v')\ne \emptyset\ \text{or}\ SN(y)\cap N(v')\ne \emptyset.
\end{align}



%By \eqref{B2Ore} and Fact \ref{MMProperty}, we have that
%begin{align}\label{SNproperty}
%    \text{Any vertex $v\in R$ with $vv'\in M$, there is at least a vertex $z\in \{x,y\}$ such that $SN(z)\cap N(z')\ne \emptyset$.}\tag{\ensuremath{*}}
%\end{align} 
$\bullet$\ Assume  $v'\notin R$. Then $v'\in SN(x)\cup SN(y)$. Without loss of generality, assume that $v'\in SN(x)$%and $v'\notin SN(y)$
, that is, $xv\in E(G)$ and $yv\notin E(G)$. Suppose that $N(v)\cap SN(y)\ne \emptyset$. Then choose $u\in N(v)\cap SN(y)$ with $uu'\in E(M)$. Hence $M':=(M\setminus \{vv',uu'\})  \cup \{vu,u'y\}$  is also a maximum matching of $G$. Notice that $v\in N(x)\cap N(v')$. Therefore, under the matching $M'$ we have $SN(x,v')\neq \emptyset$, contradicting our initial assumption in Case 2. Thus, $N(v)\cap SN(y)= \emptyset$. Together with $vy\notin E(G)$, we know $N(v)\cap V(G_2)=\emptyset$. Thus, one may add such a vertex $v$ to $V(G_1)$. % Let us consider vertices $x,v'$, which do not belong to $V(M')$. Thus, we get $u\in SN(x)\cap SN(v')$ in matching $M'$, that is, $SN(x,v')\ne \emptyset$. It is a contradiction. 

$\bullet$\ Assume  $v'\in R$. Then $xv,yv\notin E(G)$. By a similar argument as \eqref{eitheror}, we obtain that
\begin{align}\label{either2}
    \text{either}\ SN(x)\cap N(v)\ne \emptyset\ \text{or}\ SN(y)\cap N(v)\ne \emptyset. 
\end{align}
We claim that 
%$$
%SN(x)\cap (N(v)\cup N(v'))= \emptyset \text{ or }SN(y)\cap (N(v)\cup N(v'))= \emptyset.
%$$
%therwise, one 
none of the following holds:
\begin{enumerate}
    \item[{\rm (1)}] $SN(x)\cap N(v)\ne \emptyset$ and $SN(y)\cap N(v')\neq \emptyset$;
    \item[{\rm (2)}] $SN(x)\cap N(v')\ne\emptyset$ and $SN(y)\cap N(v)\ne \emptyset$. 
\end{enumerate}
%\vspace{-0.2cm}
Otherwise, without loss of generality, suppose the former case happens. Let $a\in SN(x)\cap N(v)$ and $b\in SN(y)\cap N(v')$, and assume  $aa',bb'\in E(M)$. Then $M':=(M\setminus\{aa',bb',vv'\})\cup \{xa',av,v'b,b'y\}$ is a perfect matching of $G$, a contradiction. Together with \eqref{eitheror} and \eqref{either2}, we obtain that 
\begin{align}\label{belong}
    \text{either}\ SN(x)\cap (N(v)\cup N(v'))= \emptyset\ \text{or}\ SN(y)\cap (N(v)\cup N(v'))= \emptyset.
\end{align}
Together with $xv,xv',yv,yv'\notin E(G)$, we have either $(N(v)\cup N(v'))\cap V(G_1)=\emptyset$ or $(N(v)\cup N(v'))\cap V(G_2)=\emptyset$. Therefore, $v$ and $v'$ can be assigned to $G_1$ or $G_2$ based on where their neighbors are located, which preserves the absence of edges between $G_1$ or $G_2$. % based on their neighborhood adjacencies. %Thus, $v,v'$ can be added to $G_1$ or $G_2$ according to their neighbors. 

By sequentially applying the above operation to each vertex in 
$R$, we can assign them to either 
$G_1$ or $G_2$ accordingly. In order to preserve the disconnectivity  between $G_1$ and $G_2$, we only need to show that in the above process, for each edge in $G[R]$, %$uv\in E(G[R])$, $u$ and $v$ 
its two endpoints are assigned to the same part. Let $uv\in E(G[R])$ with $uu',vv'\in E(M)$. We proceed our proof by considering the following three subcases.

%Thus, in order to finish the proof, it suffices to prove that 

%In fact, if the vertex set $R$ forms an independent set in $G$, then $G$ can be partitioned into two connected odd components, as desired. So, in what follows, we consider that there exists an edge $uv\in E(G[R])$, where $uu', vv'\in E(M)$. In order to preserve the disconnectivity  between $G_1$ and $G_2$, we will  show that in the above process, $u$ and $v$ are assigned to the same part.
%$$
%\text{either}\ \big(\bigcup_{w\in \{u,u',v,v'\}}N(w)\big)\cap V(G_1)=\emptyset\ \text{or}\ \big(\bigcup_{w\in \{u,u',v,v'\}}N(w)\big)\cap V(G_2)=\emptyset.
%$$
%. In order to put $u$ and $v$ in suitable parts, it suffices to show that 
%$uu'$ and $vv'$ are not adjacent to vertices in the same part. 


\medskip
{\bf Subcase 2.1.} $u',v'\notin R$.
\medskip

 %we obtain that $N(u)\cap SN(y)=\emptyset$ if $u'\in SN(x)$, and $N(u)\cap SN(x)=\emptyset$ if $u'\in SN(y)$. %Recall that in the above argument, we assumed  $v'\in SN(x)$ and $N(v)\cap SN(y)=\emptyset$. %Hence $vy\notin E(G)$ and $N(v)\setminus \{x\}\subseteq SN(x)\cup R$. 
We claim that there exists a vertex $z\in \{x,y\}$ such that 
$
u',v'\in SN(z).
$ 
Otherwise, suppose that $u'\in SN(x)$ and $v'\in SN(y)$. By the argument in the above,  we obtain that $N(u)\cap SN(y)=\emptyset$ and $N(v)\cap SN(x)=\emptyset$. %, which implies that $uy\in E(G)$ and $ux \notin E(G)$. %We will find a matching $M'$ such that $SN(y)\cap SN(w)\ne \emptyset$ for some vertex $w\in \{u',v'\}$. 
Without loss of generality, assume that $d(x)\le d(y)$. Then $d(x)\le \frac{n}{2}$. It follows from $v\in R$ that $xv'\notin E(G)$. Together with  \eqref{B2Ore} one has $d(v')\geq \frac{(1-2\gamma)n}{2}$. By $v'\in SN(y)$ and Fact~\ref{MMProperty}, we obtain $N(v')\cap SN(x)=\emptyset$. Thus,  $N(v')\subseteq SN(y)\cup R$. Notice that $|R|\leq \gamma n$.  Hence, there exists a vertex $a\in SN(y)\cap N(v')$. Therefore, $M':=(M\setminus \{aa',uu',vv'\})\cup \{uv,v'a,a'y\}$ is also a maximum matching in $G$. But $SN(x,u')\neq \emptyset$, a contradiction. This implies that there exists a vertex $z\in \{x,y\}$ such that 
$
u',v'\in SN(z). 
$
%which further implies $N(v)\cap SN(w)=\emptyset$ and $N(u)\cap SN(w)=\emptyset$, where $w\in \{x,y\}\setminus \{z\}$. 
Therefore, $u$ and $v$ are assigned to the same part by using the above argument. 

%ince $uy\in E(G)$ and $ux \notin E(G)$, then $d(u')\ge \frac{(1-2\gamma)n}{2}$. Since $N(u')\subseteq SN(x)\cup R$ and $|R|\le \gamma n$, then there exists $a'\in SN(x)\cap N(u)$. The matching $M':=(M\setminus \{aa',uu',vv'\})\cap \{uv,u'a,a'x\}$ also is maximum, where $aa'\in M$. The vertices $v',y$ are the vertices we need.

\medskip
{\bf Subcase 2.2.} $u'\in R,\,v'\notin R$.
\medskip

By a similar discussion as Subcase 2.1, we can show that $u$ and $v$ are assigned to the same part. % by using the above argument. $(N(u)\cup N(u'))\cap SN(y)=\emptyset$. %Otherwise, by \eqref{belong} we have $(N(u)\cup N(u'))\cap SN(x)=\emptyset$. 

\medskip
{\bf Subcase 2.3.} $u',v'\in R$.
\medskip

We claim that there exists a vertex $z\in \{x,y\}$ such that 
$$
SN(z)\cap (N(v)\cup N(v'))=\emptyset\ \text{and}\ SN(z)\cap (N(u)\cup N(u'))=\emptyset.
$$ 
Otherwise, based on \eqref{belong}, one may assume that $SN(x)\cap (N(v)\cup N(v'))=\emptyset$ and $SN(y)\cap (N(u)\cup N(u'))=\emptyset$. Furthermore, $SN(y)\cap (N(v)\cup N(v'))\neq \emptyset$ and $SN(x)\cap (N(u)\cup N(u'))\neq \emptyset$. By $uv\in E(G)$ and the maximality of $M$, one has either $SN(y)\cap N(v')=\emptyset$ or $SN(x)\cap N(u')=\emptyset$. 
Assume, without loss of generality, that $SN(y)\cap N(v')=\emptyset$. %Then $SN(y)\cap N(v)\neq \emptyset$ and  $N(v')\subseteq SN(x)\cup R$. 
Together with $SN(x)\cap N(v')=\emptyset$, one has $N(v')\subseteq R$. Notice that  $xv',yv'\notin E(G)$. By  \eqref{B2Ore}, one has $$d(x),d(y)\geq n-2\gamma n,$$
which contradicts \eqref{degreesum}. % If $d(y)\geq d(x)$, then by \eqref{B2Ore} we have $d(v')\geq \frac{(1-2\gamma)n}{2}$. Hence $N(v')\cap SN(x)\neq \emptyset$, a contradiction. Therefore, $d(x)>d(y)$.  
%Choose $a\in SN(y)\cap N(v)$ with $aa'\in E(M)$, and choose $b\in SN(x)\cap (N(u)\cup N(u'))$ with $bb'\in E(M)$.  If $b\in SN(x)\cap N(u')$, then let $M':=(M\setminus \{aa',uu',bb'\})\cup \{xb,b'u',ya'\}$. Notice that under the matching $M'$ we have $SN(u,a)\neq \emptyset$, a contradiction. Hence $SN(x)\cap N(u')=\emptyset$, which implies  $N(u')\subseteq SN(y)\cup R$. Recall that $d(x)\geq d(y)$By Ore condition, we have $d(u')\geq \frac{(1-2\gamma)n}{2}$. Hence $N(u')\cap SN(y)\neq \emptyset$, a contradiction. 

With the above discussion, by assigning vertices in $R$ to $G_1$ or $G_2$ accordingly, we know that $G$ can be decomposed into two connected components, say $C_1$ and $C_2$, each of which has odd order.
%there exists a vertex $z\in \{x,y\}$ such that 
%$$
%SN(z)\cap (\cup_{w\in \{v,v'\}\cap R}N(w))=\emptyset\ \text{and}\ SN(z)\cap (\cup_{w\in \{u,u'\}\cap R}N(w))=\emptyset. 
%$$
%Hence by adding adjacent vertices in the same part, we get two connected components of $G$, each of which has odd order. 
Furthermore, by Ore's condition, if $|C_i|\leq \frac{1-\gamma}{2}n$ for some $i\in [2]$, then $C_i$ is a clique, as desired.
\end{proof}

\end{document}
